\section{Відповіді на контрольні запитання}

\begin{description}
    \item[1. Що таке змінна в мові Java та яку інформацію визначає її тип.] ~\\
    Змінна — це іменована область пам'яті, призначена для зберігання певного значення під час виконання програми. Тип змінної визначає множину значень, які вона може набувати, обсяг виділеної пам'яті, а також набір операцій, які дозволено виконувати над цим значенням.

    \item[2. Чим відрізняються примітивні типи від об'єктних типів.] ~\\
    Примітивні типи безпосередньо зберігають свої прості значення у стеку (\textit{Stack}) і не є об'єктами. Об'єктні (посилальні) типи зберігають лише посилання (адресу) на об'єкт, який фізично розміщується в динамічній пам'яті — кучі (\textit{Heap}). Крім того, об'єктні типи можуть містити методи та набувати значення \texttt{null}.

    \item[3. Які вісім типів даних у Java є примітивними.]
    У мові Java існує рівно вісім типів даних у Java є примітивними:
    \begin{itemize}
        \item Цілочисельні: \texttt{byte}, \texttt{short}, \texttt{int}, \texttt{long};
        \item Дробові (з плаваючою комою): \texttt{float}, \texttt{double};
        \item Символьний: \texttt{char};
        \item Логічний: \texttt{boolean}.
    \end{itemize}

    \item[4. Для чого використовується тип \texttt{int}.] ~\\
    Тип \texttt{int} використовується для представлення 32-бітних цілих чисел зі знаком. Він є стандартом за замовчуванням у Java для більшості звичайних математичних обчислень, лічильників циклів та збереження дискретних кількостей об'єктів.

    \item[5. У чому полягає відмінність між \texttt{float} та \texttt{double}.] ~\\
    Тип \texttt{float} є типом одинарної точності, займає 32 біти в пам'яті й потребує суфікса \texttt{f} при оголошенні літералів. Тип \texttt{double} є типом подвійної точності, займає 64 біти, забезпечує значно вищу точність обчислень і є стандартом за замовчуванням для дробових чисел у Java.

    \item[6. Для чого використовується тип \texttt{boolean}.] ~\\
    Логічний тип \texttt{boolean} призначений виключно для зберігання результатів умовних виразів та логічних операцій. Він може набувати лише одного з двох можливих значень: \texttt{true} (істина) або \texttt{false} (хибність).

    \item[7. Чим відрізняється \texttt{char} від \texttt{String}.] ~\\
    Тип \texttt{char} є примітивним, займає 16 біт і зберігає строго один символ Unicode в одинарних лапках (наприклад, \texttt{'A'}). Клас \texttt{String} є об'єктним типом даних, оперує послідовністю символів у подвійних лапках (наприклад, \texttt{"Java"}) та інкапсулює методи для роботи з текстом.

    \item[8. Чому рядок є об'єктним типом.] ~\\
    Рядок у Java є екземпляром повноцінного класу \texttt{String}, а не просто масивом примітивів чи базовим типом. Рядки мають власні методи обробки, підтримують механізм незмінності (\textit{Immutability}) та зберігаються у кучі, де оптимізуються спеціальним пулом рядків (\textit{String Pool}).

    \item[9. Що таке класи-обгортки.] ~\\
    Класи-обгортки (\textit{Wrapper Classes}) — це спеціальні класи з пакету \texttt{java.lang}, які представляють примітивні типи даних у вигляді повноцінних об'єктів. Кожному примітиву відповідає свій об'єктний аналог (наприклад, \texttt{int} $\rightarrow$ \texttt{Integer}, \texttt{double} $\rightarrow$ \texttt{Double}), що дозволяє використовувати їх там, де потрібні виключно об'єкти.

    \item[10. Що таке автоупаковування та автораозпаковування.] ~\\
    \textbf{Автоупаковування} (\textit{Autoboxing}) — це автоматичне перетворення примітивного типу в об'єкт відповідного класу-обгортки (наприклад, присвоєння значення \texttt{int} змінній типу \texttt{Integer}). \\
    \textbf{Авторозпаковування} (\textit{Auto-unboxing}) — зворотний процес, при якому об'єкт класу-обгортки автоматично трансформується у примітивне значення при виконанні арифметичних операцій.
    \item[11. Що відбувається під час перетворення \texttt{double} у \texttt{int}.] ~\\
    Під час явного приведення типу \texttt{double} у \texttt{int} відбувається примусове відсікання (відкидання) всієї дробової частини числа, а не його математичне округлення. Значення просто урізається до цілого числа, що може призвести до втрати точності даних (наприклад, \texttt{(int) 25.75} перетвориться на \texttt{25}).

    \item[12. Для чого використовується ключове слово \texttt{final}.] ~\\
    Ключове слово \texttt{final} використовується для створення констант, тобто значень, які не можуть бути змінені після їхньої першої ініціалізації. Оголошення змінної або поля класу з модифікатором \texttt{final} блокує будь-які спроби повторного переприсвоєння значення на етапі компіляції.

    \item[13. Чим відрізняється локальна змінна від поля класу.] ~\\
    Локальна змінна оголошується безпосередньо всередині методу або блоку коду, існує лише під час його виконання і зберігається у стеку. Вона не має значень за замовчуванням і вимагає обов'язкової ручної ініціалізації перед використанням. Поле класу оголошується всередині тіла класу, визначає стан об'єкта у кучі й автоматично ініціалізується значеннями за замовчуванням при створенні об'єкта.

    \item[14. Що таке конструктор класу.] ~\\
    Конструктор — це спеціальний метод класу, який автоматично викликається оператором \texttt{new} під час створення нового об'єкта для ініціалізації його початкового стану. Конструктор не має типу повертаємого значення, а його назва повинна строго збігатися з ім'ям самого класу.

    \item[15. Для чого використовується ключове слово \texttt{this}.] ~\\
    Ключове слово \texttt{this} є прихованим посиланням на поточний екземпляр об'єкта, всередині якого викликається метод або конструктор. Найчастіше воно застосовується для розв'язання конфлікту імен, коли назви вхідних аргументів конструктора повністю збігаються з іменами полів класу (наприклад, \texttt{this.name = name}).

    \item[16. Що таке об'єкт.] ~\\
    Об'єкт — це конкретний будівельний блок об'єктно-орієнтованої програми, який створюється в динамічній пам'яті (кучі), є реальним втіленням (екземпляром) свого абстрактного класу та має власний унікальний стан і визначену поведінку.

    \item[17. Чим відрізняється примітивна змінна від змінної об'єктного типу.] ~\\
    Примітивна змінна безпосередньо містить у собі конкретне просте значення (число, символ або логічний прапорець). Змінна об'єктного типу містить у собі лише адресу (посилання), яка вказує на місце розробки реального об'єкта у кучі. При копіюванні примітиву копіюється саме значення, а при копіюванні об'єктної змінної — лише адреса посилання на той самий об'єкт.

    \item[18. Що означає значення \texttt{null}.] ~\\
    Значення \texttt{null} — це спеціальний літерал, який може набувати будь-яка змінна об'єктного (посилального) типу. Воно офіційно означає, що дана змінна наразі є порожньою і не посилається на жоден реальний об'єкт у кучі.

    \item[19.] \textbf{Чому звернення до методу через посилання зі значенням \texttt{null} може спричинити \texttt{NullPointerException}.} ~\\
    Оскільки значення \texttt{null} вказує на відсутність об'єкта, спроба викликати через таке порожнє посилання будь-який метод або звернутися до поля змушує віртуальну машину Java шукати дані за несучуючою адресою. Не маючи фізичного об'єкта для виконання операції, JVM негайно зупиняє роботу програми та генерує критичну помилку \texttt{NullPointerException}.

    \item[20. Що таке масив.] ~\\
    Масив у мові Java — це повноцінний об'єктний тип даних, який є впорядкованою структурою для зберігання фіксованої кількості елементів одного типу. Розмір масиву жорстко визначається під час його створення, не може змінюватися в процесі роботи програми, а доступ до його елементів здійснюється за допомогою індексів.
    
    \item[21. Чому індексація масивів у Java починається з нуля.] ~\\
    Індексація з нуля базується на низькорівневій логіці обчислення адрес пам'яті процесором. Індекс елемента масиву насправді є коефіцієнтом \textbf{зміщення (offset)} відносно початкової адреси першого елемента. Оскільки перший елемент знаходиться безпосередньо на початку виділеного блоку пам'яті, його зміщення дорівнює нулю.

    \item[22. Чим масив об'єктів відрізняється від масиву примітивних значень.] ~\\
    Масив примітивних значень безпосередньо містить усередині своїх комірок самі прості значення (наприклад, числа типу \texttt{int}). Масив об'єктів містить у своїх комірках не самі об'єкти, а лише \textbf{адреси посилань} на них. Самі об'єкти при цьому створюються та зберігаються окремо в кучі (\textit{Heap}), а за замовчуванням комірки масиву об'єктів ініціалізуються значенням \texttt{null}.

    \item[23. Що таке інкапсуляція.] ~\\
    Інкапсуляція — це один із трьох фундаментальних принципів ООП, який полягає в логічному об'єднанні даних (полів стану) та методів для їхньої обробки всередині однієї сутності (класу), а також у прихованні внутрішньої реалізації об'єкта від прямого неконтрольованого доступу з боку зовнішніх компонентів програми.

    \item[24. Для чого поля класу оголошують як \texttt{private}.] ~\\
    Оголошення полів як \texttt{private} є інструментом забезпечення інкапсуляції. Це повністю перекриває можливість прямої некоректної модифікації внутрішнього стану об'єкта ззовні (наприклад, встановлення негативного віку або ціни). Контроль доступу та валідація даних у такому випадку примусово здійснюються через відкриті публічні методи інтерфейсу.

    \item[25. Що таке наслідування.] ~\\
    Наслідування — це принцип ООП, який дозволяє створити новий (дочірній) клас на основі вже існуючого (базового) класу. Дочірній клас при цьому автоматично переймає структуру полів та методи батьківського класу, отримуючи можливість розширювати їх або перевизначати їхню поведінку, що забезпечує повторне використання коду.

    \item[26. Для чого використовується ключове слово \texttt{extends}.] ~\\
    Ключове слово \texttt{extends} використовується в оголошенні класу для явного зазначення того, що цей клас є дочірнім (підкласом) і успадковує функціональність іншого конкретного суперкласу (наприклад, \texttt{public class Student extends Person}).

    \item[27. Що таке поліморфізм.] ~\\
    Поліморфізм — це принцип ООП, який дозволяє використовувати одне й те саме посилання базового типу для взаємодії з об'єктами різних дочірніх класів. Завдяки механізму пізнього зв'язування, віртуальна машина Java (JVM) під час виконання програми автоматично визначає фактичний тип об'єкта і викликає його власну перевизначену версію методу.

    \item[28. Яку роль відіграє конструктор дочірнього класу та ключове слово \texttt{super}.] ~\\
    Конструктор дочірнього класу відповідає за повну ініціалізацію об'єкта. Ключове слово \texttt{super} використовується для виклику конструктора батьківського класу (наприклад, \texttt{super(name)}) і обов'язково повинно стояти першим рядком. Це гарантує, що спочатку в пам'яті будуть коректно створені та ініціалізовані базові поля суперкласу, і лише потім виконається специфічна логіка дочірнього компонента.
\end{description}

\end{description}
